\chapter{كيف تتعلم الآلة الحركة}
% full version: chapters/15_motorlearning.tex

رأينا طريقتين للتعلم. بلا معلم، ”معًا يعني مترابطين“، تحفظ الآلة ما يتكرر كثيرًا. ومع الدوبامين تتعلم من المفاجآت: جاءت النتيجة أفضل أو أسوأ مما كان متوقعًا. لكن حين تخطئ وأنت ترمي الكرة في السلة، لا تعرف فقط أن النتيجة سيئة. تعرف \emph{بكم} أخطأت و\emph{في أي اتجاه}. هذا هو المعلم الثالث: خطأ له اتجاه، خاص بكل مخرج.

\section*{سلك الخطأ}

في الدماغ دائرة منفصلة لتعلم الحركات، وبنيتها غير مألوفة إطلاقًا. عدد هائل من عصبونات \nt{GLUT} الصغيرة، عشرات المليارات، يفرد الإشارات الواردة في قائمة طويلة جدًا من السمات. وعصبونات الخرج في هذه الدائرة كابحة، من نوع \nt{GABA}، وعددها عشرات الملايين، وكل منها يصغي إلى مئة ألف مدخل. وإلى كل عصبون خرج يمتد سلك خاص واحد بالضبط: سلك الخطأ. إذا وصل الخطأ والمدخل نشط، ضعف هذا المدخل. هذه هي القاعدة الكلاسيكية ”قلّل الخطأ“ التي يعلّم بها المهندسون المرشحات التكيفية.

\pencilfig{15}{\pencilart{15}}{الإخفاق لا يقول ”سيئ“ فحسب، بل ”بكم وفي أي اتجاه“.}

هذا المعلم أسرع من معلم الدوبامين: لكل مخرج خطؤه، ولا حاجة إلى استخلاص الحصة الخاصة من إشارة مشتركة للجميع. لكن الخطأ يصل متأخرًا، فتحتاج الوصلة مرة أخرى إلى ”ملصق“ يتذكر أنها شاركت. وتبين التجارب أن مدة هذا الملصق مضبوطة على تأخر الخطأ.

\section*{النموذج الداخلي}

ماذا تتعلم دائرة كهذه؟ نموذجًا للجسم نفسه: أي نتيجة سيعطيها الأمر. ومعه لا حاجة إلى انتظار المستشعرات المتأخرة: يمكن التنبؤ بالنتيجة مسبقًا. ووفق إحدى الفرضيات، لهذا لا تستطيع أن تدغدغ نفسك: الدماغ يعرف مسبقًا ما سيحدث، فيطرح المتنبَّأ به.

\section*{متعلم سريع ومتعلم بطيء}

البس نظارة تزيح الصورة جانبًا، وارمِ السهام. الرميات الأولى تنحرف، لكنك تتكيف خلال عشرات المحاولات. انزع النظارة فتعود الإخفاقات، في الاتجاه الآخر هذه المرة. ويبدأ الأمر المثير بعد ذلك. التجارب يصفها جيدًا متعلمان معًا: السريع يتعلم سريعًا وينسى سريعًا، والبطيء يتعلم ببطء لكنه يتذكر طويلًا. إذا أعدت التعلم بالعكس لمدة قصيرة بعد تكيف طويل، عاد التصحيح الكلي إلى الصفر، لكن المهارة القديمة تطفو بعد ذلك جزئيًا من تلقاء نفسها: السريع نسي، والبطيء يتذكر. النموذج ذو المتعلم الواحد لا يقدر على هذا، أما البشر فيقدرون.

\pencilfig{115}{%
% figures/tworate_*.dat from models/motor_learning.py: adapt 200 throws, reverse until the sum is back to zero, then no errors at all
\begin{scope}[x=0.026cm, y=1.45cm]
\fill[pcGray, opacity=0.08] (220,-1.1) rectangle (226,1.1);
\draw[pen] (0,0) -- (340,0);
\draw[pcGray, line width=0.4pt] plot file {figures/tworate_p.dat};
\draw[penlite=pcAmber] plot file {figures/tworate_fast.dat};
\draw[penlite=pcBlue] plot file {figures/tworate_slow.dat};
\draw[pcGray!40!black, line width=1.1pt] plot file {figures/tworate_net.dat};
\draw[pcGray, densely dashed, line width=0.4pt] (226,-1.1) -- (226,1.1);
\end{scope}
\node[lbl, anchor=south] at (3.1,1.47) {النظارة على العينين};
\node[lbl, anchor=north west, align=left] at (6.3,-0.55) {إعادة تعلم\\قصيرة};
\node[lbl, anchor=south west] at (6.0,1.47) {لا أخطاء};
\node[lbl, text=pcBlue, anchor=north] at (4.4,0.8) {البطيء};
\node[lbl, text=pcAmber!80!black, anchor=south west] at (1.15,0.5) {السريع};
\node[lbl, text=pcGray!40!black, anchor=south] at (7.7,0.95) {المجموع طفا};
}{تصحيح المتعلمَين في النموذج. بعد إعادة التعلم القصيرة يكون المجموع عند الصفر، لكن البطيء يتذكر القديم، فيطفو من تلقاء نفسه.}

لماذا متعلمان؟ وفق الفرضية، إنها وصفة الملاح نفسها من الفصل التاسع: العالم يتغير بسرعة وببطء معًا، ومن الحكمة تتبع الاثنين.

\fullversion{15}
